\documentclass[10pt, a5paper, twoside]{ltjsarticle}

%\usepackage{okumacro}
\usepackage{enumitem}
\usepackage{amssymb}
\usepackage{amsmath}
\usepackage{amsfonts}
\usepackage{amsthm}
\usepackage{mathrsfs}
\usepackage{bm}
\usepackage{xurl}
\usepackage{here}
\usepackage{graphicx}
\usepackage{wrapfig}
\usepackage{makeidx}
\usepackage{braket}
%\usepackage{ascmac}
\usepackage{fancyhdr}
\usepackage{tcolorbox}
\usepackage{tikz}
%\usepackage{otf}
\usepackage{luatexja-otf}
\usepackage{luatexja-ruby}
\usepackage[top=20truemm,bottom=20truemm,left=15truemm,right=15truemm]{geometry}
\usepackage{mathrsfs}

\pagestyle{fancy}
	\fancyhead{}
	\fancyhead[RE]{円城塔研究}
	\fancyhead[LO]{資料：コンメンタール「Goldberg Invariant」}
	\fancyhead[LE, RO]{\thepage}
	\fancyfoot{}
	\fancyfoot[LE, RO]{\footnotesize{The Journal of EnJoeToh Research, Vol.4, No.1, 2026} }

\tcbuselibrary{breakable, theorems, skins}

\newtcbtheorem{axm}{公理}{enhanced, breakable=true}{axm}
\newtcbtheorem[use counter from=axm]{dfn}{定義}{enhanced, breakable=true}{dfn}
\newtcbtheorem[use counter from=axm]{thm}{定理}{enhanced, breakable=true}{thm}
\newtcbtheorem[use counter from=axm]{sys}{系}{enhanced, breakable=true}{sys}
\newtcbtheorem[use counter from=axm]{hyp}{予想}{enhanced, breakable=true}{hyp}
\newtcbtheorem[use counter from=axm]{ths}{提唱}{enhanced, breakable=true}{ths}
\newtcbtheorem[use counter from=axm]{prn}{原理}{enhanced, breakable=true}{prn}
\newtcbtheorem[use counter from=axm]{clm}{主張}{enhanced, breakable=true}{clm}
\newtcbtheorem[use counter from=axm]{prb}{問題}{enhanced, breakable=true}{prb}
\newtcbtheorem[use counter from=axm]{prf}{証明}{enhanced, breakable=true, colback=white, attach boxed title to top left={xshift=2mm, yshift=-2mm}}{prf}

\newcommand{\hyphen}{\insertkanjiskip--\insertkanjiskip}

\newcommand{\cross}[1]{%
\mbox{\vbox{\kern 1pt\hbox{\vbox{\hrule
\kern 2pt\hbox{\ensuremath{\vphantom{b}#1}\kern 2pt}}\vrule\kern 1pt}}}\,}

\begin{document}

	~ %強制改行

	\begin{center}

		\Large{コンメンタール「Goldberg Invariant」}

		\vspace{3mm}

		\large{Commentary of short story \textit{Goldberg Invariant}}

		\vspace{3mm}

		\large{下村思游}

	\end{center}

	\vspace{3mm}

	\section{解説}

        頁数・行数は『Boy's Surface』紙版\cite{boys}に準拠した．

        \subsection{題名}

            自作解題「What is the Name of This Rose?」にもある通り，バッハの『ゴルトベルク変奏曲』(Goldberg Variations)のVariationsをInvariant(不変量)とした駄洒落に由来する．

            『ゴルトベルク変奏曲』は，厳格で純理論的ではあったが，人前では滅多に演奏されることのない曲という評価だった\cite{gould}が，グレン・グールドによる演奏が世界的な大ヒットとなった．また，グールドのキャリアの最初を飾ったデビューアルバムがこの『ゴルトベルク変奏曲』であり，結果的に最後のアルバムとなったのも『ゴルトベルク変奏曲』であることは，偶然ではあるが非常に興味深い．さらに，グールド本人のキャリア自体も『ゴルトベルク変奏曲』の進行\footnote{「主題のアリアで始まり，各種の変奏を経て，ダ・カーポのアリアが回帰して終わる」\cite{miyazawa}p.133}に類似するとの指摘がある\cite{miyazawa}．また，初期円城塔に頻出する天才表象という点からは，グールドが様々な奇行で知られた人物であることにも注意したい．

        \subsection{エピグラフ}

            \begin{quote}
                「美都久幾能〜」
            \end{quote}

            江戸時代末期の越後の僧侶，良寛の短歌．全文は「みづくきの跡も涙にかすみけり在りし昔のことを想ひて」．

            これは父\footnote{俳人としても著名で，小林一茶とも交流があった．}の句「朝霧に一段ひくし合歓の花　以南」の左隅に小さく添え書きされていたものである\cite{togo}．“みづくき”とは筆跡のこと．良寛は，父の遺した書き物を見ながら，生前の父を思い出している\cite{ryo100}．良寛は名主の長男として生まれ，裕福だったのだが，若い頃から愚鈍と放蕩で知られ，しかも突如出家するなどの親不孝を重ねた．父の筆跡を前に，自身のかけた迷惑に思いを馳せたのだろう．

            エピグラフとしては，霧島との過去を懐かしく思っていることを象徴するか．ただし，最終盤で明かされる通り，“わたし”は霧島との敵対を望んでおり，良寛の短歌とは風情が異なる．

        \subsection{p.78, l.2}

            \begin{quote}
                「霧島梧桐」
            \end{quote}

            綴りはGodoh Kirishimaであり，本作の物語は，姿を消した霧島を待つ，という構成となっている．ここから，サミュエル・ベケットによる戯曲『ゴドーを待ちながら』が由来であることが強く示唆される．また漢字は歌人，河東碧梧桐からの借用か．さらに，計算機というモチーフからは物理学者・計算機科学者の後藤英一や，計算機科学者エドガー・ダイクストラによる有名な論文``Go To Statement considered harmful''（「Goto文は有害であるように思われる」）\cite{dijkstra}も連想されるが，これらはやや飛躍の気配が強い．

        \subsection{p.78, l.2}

            \begin{quote}
                「GRAPE64」
            \end{quote}

            詳細は\ref{grape}節で説明するが，これは東京大学大学院総合文化研究科で開発された，重力多体系専用計算機であるGRAPEシリーズに由来する．

        \subsection{p.78, l.5-8}

            \begin{quote}
                「戦役の推移の概観については，〜」
            \end{quote}

            ある分野の創始者がジャーナルでその分野の特集を監修し，執筆したレビュー記事がその分野の教科書代わりとなって大学院生の必読文献になる，というのは少なくとも物理学や数学ではよく見られる光景である．

            なお，作中に登場するReview of Modern Physicsは実在する雑誌．アメリカ物理学会(American Physical Society, APS)が発行する査読誌で，その名の通りレビュー論文誌．Physical Review Letters, Physical Review A-Eに並んで最も重要な物理学誌のひとつ．

        \subsection{p.78, l.9}

            \begin{quote}
                「サンクトペテルブルク線」
            \end{quote}

            第二次大戦の独ソ戦における激戦地，レニングラード（当時）からか．

            のちほど，人類の戦線として環太平洋線，五大湖線も挙げられている．いずれも水域に面しているが，水域と代替数学との間の関係性は不明．

        \subsection{p.78, l.10}

            \begin{quote}
                「ゴルトベルクの名を冠する定理の数は，おおよそ一年半に二倍の割合で増加を続けているという事実」
            \end{quote}

            経過時間に対する半導体の性能の向上率を指摘した経験則，ムーアの法則をもじったもの．

            \begin{dfn}{ムーアの法則}{moore}
                集積回路上のトランジスタの数は18ヶ月ごとに2倍になる．
            \end{dfn}

            「ムーアの限界」というのは，ムーアの法則には量子力学や熱力学による適用限界があることを指すか．本作では，これを転じて，ゴルトベルクの名を冠する定理の数の増加率はいつか鈍るはずだという意味で用いている．ムーアの法則は，後述のマーゴラス\hyphen レヴィティンの定理のほか，ランダウアーの原理によっても制限される．

        \subsection{p.78, l.17}

            \begin{quote}
                「\ltjruby{P}{ポリノミアル}から\ltjruby{NP}{ノンポリノミアル}領域へ」
            \end{quote}

            文脈から，計算複雑性理論におけるクラスであるPとNPに由来する．ここで，クラスP，NPの定義は以下の通り\cite{sip3}．

            \begin{dfn}{クラスP}{}
				Pは決定性単一テープチューリング機械によって多項式時間で判定出来る言語のクラスである．
				$$\mathrm{P} = \bigcup_{k} \mathrm{TIME}(n^k)$$
			\end{dfn}

            \begin{dfn}{クラスNP}{}
                NPは多項式時間検証装置をもつ言語のクラスである．
            \end{dfn}

            ここで，クラスPが多項式時間\footnote{多項式時間とは，時間を計算のステップ数で表したとき，多項式で表される時間のことである．例えば，判定に$n, n^2, n^3$回の計算が必要ならば，多項式時間で判定出来るが，$2^n$回の計算が必要ならば，多項式時間では判定出来ない（指数時間で判定可能）．}(polynomial time)に由来するのに対し，クラスNPは非決定性多項式時間(non-deterministic polynomial time)に由来することに注意せよ．よくある誤解として，クラスPの補集合をクラスNPと考えるものがある．この場合，未解決問題として有名なP $\neq$ NP問題は自明．

            \begin{prb}{P $\neq$ NP問題}{}
                クラスPはクラスNPに一致するか，あるいはクラスPはクラスNPの真部分集合であるか．
            \end{prb}

            P $\neq$ NP問題は，アメリカのクレイ数学研究所によって100万ドルの懸賞金がかけられたミレニアム懸賞問題のひとつである．クラスPがクラスNPの部分集合であることは自明だが，一致するのか真部分集合であるのかという点は未解決となっている．なお，多くの数学者は，クラスPはクラスNPの真部分集合であると考えている．

            一方で，作中では，Pに“ポリノミアル”，NPに“ノンポリノミアル”というルビが振られている．これは先述の誤解を踏まえたギャグであると考えられる．仮にクラスPの補集合がクラスNPであるとすると，先述の通り，P $\neq$ NP問題が直ちに解決される．これは明らかに愚かな帰結であり，極めて滑稽である．

            あるいは，NPは“多項式時間で解くことができるという証拠が存在する”というクラスなので，この延長であるNNPは“多項式時間で解くことができるという証拠が存在するという証拠が存在する”というクラスを意味しているとも取れる．これは直ちに無限後退に陥ることから，円城塔作品の読解として説得的である．

            先の仮定の文学的意図について考察する．通常の数学，例えばZFC公理系を採用した数学では，補集合の補集合が元々の集合に一致することは自明\footnote{\because \, 排中律．}．したがって，通常の数学ではクラスNNPなるものは意味不明な上，仮定したとしても自明な対象でありつまらない．逆に，クラスPに対する補集合の補集合であるクラスNNPが非自明であるなら，これは作品内の数学が排中律を採用していないことを示唆する．本作は数学と相異なる別の数学（代替数学）が対立し，強大な計算能力によって互いを塗り替え合うという物語\footnote{この着想は，のちの『エピローグ』「リアルタイムラジオ」の主題となる．}が語られる．これは明らかにグレッグ・イーガン「ルミナス」を踏襲しているが，「ルミナス」が“不安定な排中律”を採用して真偽の書き替えを描いたのに対して，（本作が排中律を採用していないことを認めるならば）本作は真偽のさらに外にある形而上的な領域での認知不能な戦闘を描くという違いがある．また，「ルミナス」がその書き換えが人知れず行われるのに対して，本作では一般社会に影響を与えているという違いが指摘される．

            さらに別の先行作としてボルヘス「トレーン，ウクバール，オルビス・テルティウス」が指摘される．ただし，「トレーン」が人文的な知識に対する攻撃だったのに対して，本作は数理的な真理に対する攻撃である点が大きく異なる．

            また，本作は，数学の定理に関する真偽値が塗り替わるダイナミクスを，その内部から描写した作品と考えることが出来る．真偽値はバイナリであり，バイナリを一定のルールで塗り替えていくダイナミクスとしてラングトンの蟻が連想される．この見立ては魅力的だが，本作がラングトンの蟻をモチーフにしているのであれば，ラングトンの蟻において現れる特徴的な秩序構造についての描写が作中に見られるはずである．残念ながら本作にはその秩序構造の描写が見られないので，解釈としては棄却すべきだろう．

        \subsection{p.79, l.2-4}

            \begin{quote}
                「警告を発する者の多くが，既にその戦闘展開規模と速度に追随できなくなった退役研究者であることは皮肉を伴う．」
            \end{quote}

            プランクの原理として知られる箴言を踏まえた表現か．ここで，プランクの原理は以下のような形で知られている．

            \begin{prn}{プランクの原理}{planck}
                科学は葬式のたびに進歩する．
            \end{prn}

            これは量子論の創始者のひとりである物理学者，マックス・プランクのものとされている．実のところ，科学にまつわる逸話の多くが変形されたものであり，このプランクの原理もプランク本人の発言の趣旨を要約したものである．プランクが実際に述べたものは以下の通り\cite{planck}．

            \begin{quote}
                新たな科学的真理が普及するのは，反対者が説得され，誤りを認めることによってではなく，反対者が徐々に死に絶えていき，新たな世代が最初からその真理を受け入れて喧伝することによってである．
            \end{quote}

        \subsection{p.80, l.1}

            \begin{quote}
                「セイタカアワダチソウじみた戦略」
            \end{quote}

            侵略的外来種として知られるセイタカアワダチソウは，根から毒素を土壌に放出することで他の植物の成長を阻害\footnote{このような作用のことをアレロパシーという．日本在来の植物ではアカマツが有名．}し，自身の繁茂する余地を確保するという生存戦略で知られている．なお，セイタカアワダチソウの繁殖密度が高い場合，自身の放出した毒素によって自身の成長が阻害される．

        \subsection{p.80, l.4-}

            \begin{quote}
                「当時の戦闘は非道く長閑なものだった．我々には機械支援なしでプログラムを組み立てる時間が与えられていたし，〜」
            \end{quote}

                本作の発表当時は何気ない記述だったかもしれないが，LLMのコーディング能力が人間を軽く凌駕するようになった今となっては，大変強い実感をもって納得させられる．


        \subsection{p.80, l.7-8}

            \begin{quote}
                「カスケード的に数千の定理を一気に生成して〜」
            \end{quote}

            カスケードとは，小さな複数の滝が複層的に連なったもののこと．ここから転じて，あるひとつの要素から，“子”が生成され，それらがさらに“孫”を生成させて急速に増殖する様を，物理学およびその周辺分野ではカスケードと表現することがある．

            私が知る実例では，素粒子物理学におけるものがある．ある粒子が崩壊して娘粒子を生み出し，それがさらに崩壊を繰り返す現象をカスケードといい，その崩壊粒子の集まりをカスケードシャワーという．

            また，このカスケードのイメージは，カントール集合によく似ている．カントール集合は再帰的なアルゴリズムによって生成可能であり，フラクタル構造の代表としても知られる．カントール集合の定義は以下の通り\cite{chaos1}．

            \begin{dfn}{カントール集合}{}
                単位区間$[0, 1]$に対して，間を取り除く操作を無限に繰り返して以下のように得られる極限集合を\textbf{カントール集合}という．
                \begin{align*}
                    C = [0, 1] \setminus \bigcup^{\infty}_{n=1} \bigcup^{3^{n-1}-1}_{k=0}\bigg(\frac{3k+1}{3^n}, \frac{3k+2}{3^n}\bigg)
                \end{align*}
            \end{dfn}

            カントール集合の性質として，閉集合でありしかもコンパクトである，連続体濃度を持つ，ルベーグ測度が0であるなどの特徴的な性質が知られている．また，カントール集合に含まれる数を三進数で表現すると，どの桁にも1が出現しない．

        \subsection{p.80, l.10-11}

            \begin{quote}
                「非道く単純であるが故に生じえたものは，複雑化してしまった事象の中から生み出される機縁を持てない．」
            \end{quote}

            ウルフラムによるオートマトンのクラス分け，およびウルフラムの提唱を強く意図した記述．ここで，ウルフラムによるオートマトンのクラス分けは以下の通り\cite{bio}．

            \begin{quote}
                \begin{enumerate}
                    \item 単調でつまらない「死滅」状態
                    \item 単純で安定な周期的状態
                    \item 非常に多くの安定点を持ち，さらに不安定な状態にもなるカオス的なパターン
                    \item 複雑な局所的構造で，動的に進化する．ある期間後に突然パターンを変えるような状態．
                \end{enumerate}
		    \end{quote}

            オートマトンは，他のクラスへと自然に遷移することはない．つまり，単純な構造から複雑な構造が自然に生じることはなく，逆に複雑な構造から単純な構造が自然に生じることもない．作中の記述は，このことを示しているものと思われる．

            また，ウルフラムによるクラス分けからは，ウルフラムの提唱が直ちに連想される．ウルフラムは，4番目に挙げた，十分複雑%\index{ふくざつ@複雑!じゅうぶんふくざつ@十分---}
            で“生きた”計算論的振舞いをするオートマトンは大体チューリング完全である，と主張した．このような主張をウルフラムの提唱%\index{うるふらむのていしょう@ウルフラムの提唱}
            といい，90年代から00年代の複雑系研究者の間で一種の常識として扱われていた．

            \begin{ths}{ウルフラムの提唱}{}
    			十分複雑なものはチューリング完全であり，それは生命である．
		    \end{ths}

            これらの話題は円城塔の作品，特に最初期の作品\footnote{「パリンプセストあるいは重ね書きされた八つの物語」「ムーンシャイン」「これはペンです」など．}によく見られる．円城塔にとって，複雑な構造であったり，同質なものが大量に並べられていたりするものは当然にチューリング完全な計算機であり，それは生命であって，意識をもっている．

            また，ある状態が，ある時点以降の状態からは取り出せない，遷移出来ないといった要素から，エデンの園定理や唯一の父問題といった話題が連想される．エデンの園定理とその周辺の話題については\ref{eden}節で詳細に説明する．

        \subsection{p.82, l.15}\label{write}

            \begin{quote}
                「書かれることを拒む記号列の抵抗．」
            \end{quote}

            作中ではこの着想を“滑稽である”として棄却しているが，円城塔作品には，書くことによって失われてしまうことへの執着がしばしば見られる．『Self-Reference ENGINE』のSelf-Reference ENGINEがその好例．

        \subsection{p.83, l.1-5}

            \begin{quote}
                「私は私の書く物自体を欺かなければならないらしい〜」
            \end{quote}

            対象に気づかれないように女性が潜入する，という点からは記紀神話においてヤマトタケルが女装してクマソタケルを討った逸話が連想される．

            また，ここから続く林檎という文字云々という話は，ゲシュタルト崩壊を指しているか．あるいは，単に病的な妄想癖の描写か．特に林檎が取り上げられているのは，アラン・チューリングが青酸化合物を塗布した林檎を食べて自殺したことを意識したものだろう．また，食事に関する妄想癖という点は，\ref{deathofgodel}節で触れる，最晩年のゲーデルとの関連が指摘される．

        \subsection{p.83, l.14}

            \begin{quote}
                「結局私にはこういう形でしかこの現象を伝えることができなかった．」
            \end{quote}

            本作の記述が錯綜し混乱しているように見えることについての，語り手による言い訳．

        \subsection{p.84, l.3}

            \begin{quote}
                「春の雨があがり，宇宙が晴れ上がる．」
            \end{quote}

            前半部分は一般的な語用であり，その出典は不明だが，後半部分は宇宙論における“宇宙の晴れ上がり”に由来している．

            宇宙が誕生してしばらくの間,宇宙は高エネルギー状態にあり，低エネルギー帯では結合しない電子と光子が結合してしまうことで，光が直進出来ない状態にある．宇宙の膨張とともに宇宙が冷えていき，電子と光子が脱結合(decouple) し，光が直進出来るようになり澄み渡っていく時代のことを，宇宙の晴れ上がりという．“霧”によって遠方の情景が遮られて観測出来なくなる，という作中の描写は，この宇宙の晴れ上がりを前提としていることを強く示唆する．同様の描写は，のちの「ムーンシャイン」にも見られる．

        \subsection{p.84, l.7-}

            \begin{quote}
                「私たちにはキャサリンA，キャサリンB，キャサリンCとそれぞれ名前がつけられている．〜」
            \end{quote}

            同じ姿をした同じ名前の存在が大量に存在するのは，ジョン・ヴァーリイ「バービーはなぜ殺される」との関連が指摘される．バービー(Barby)に対してキャサリン(Catherine)という字面も連続性を示唆する．

        \subsection{p.85, l.5}

            \begin{quote}
                「この番組は，キャサリンAの提供でお送りしています．」
            \end{quote}

            初期円城塔によく見られる，突然のギャグ．

        \subsection{p.85, l.7-12}

            \begin{quote}
                「原初の渾沌を凝っと睨んで，〜」
            \end{quote}

            日本書紀の巻第一神代上「古に天地未だ\ruby{剖}{わか}れず，\ruby{陰陽}{めを}分れざりしとき，\ruby{渾沌}{まろか}れたること鶏子の如くして，\ruby{溟涬}{ほのか}にして\ruby{牙}{きざし}を含めり．其れ\ruby{清陽}{すみあきらか}なるものは，\ruby{薄靡}{たなび}きて天と為り，\ruby{重濁}{おもくにご}れるものは，\ruby{淹滞}{つつ}ゐて地と為るに及びて，\ruby{精妙}{くはしくたへ}なるが合へるは\ruby{搏}{むらが}り易く，重濁れるが凝りたるは\ruby{竭}{かたま}り難し．」\cite{nihonshoki}に由来する記述．混沌という語は元々中国の様々な神話に由来し，この日本書紀の記述はそれらの中から日本神話に近いものを集めてきてここにまとめたもの．また，「凝っと睨んで」という箇所は，物理学者や数学者が数式をしばしば“ぐっと睨む”ことに由来するか．ある数式をよく観察して，物理学的・数学的洞察を得ることを，ぐっと睨むと表現する．

        \subsection{p.85, l.13-17}

            \begin{quote}
                「太初に言ありき，言は我らと供にあり，言は我らなりき．」
            \end{quote}

            文語訳新約聖書の「ヨハネ伝福音書」（ヨハネによる福音書）第1章1-3節「\ruby{太初}{はじめ}に\ruby{言}{ことば}あり，言は神と\ruby{偕}{とも}にあり，言は神なりき．この言は太初に神とともに在り，万の物これに由りて成り，成たる物に一つとして之によらで成りたるはなし．」\cite{bible}をやや改変しつつ，記紀神話の国産み神話に混ぜ込んでいったもの．意外なことに両者は字面がかなり似ており，現代神話としてのSFという演出以外にも，単純に連想しただけという可能性も指摘される．

        \subsection{p.86, l.1-7}

            \begin{quote}
                「私たちの仕事は，砂嵐を許々袁々呂々と掻き回して，〜」
            \end{quote}

            記紀神話の記述を熱力学・統計力学的に解釈し直し，熱力学・統計力学に基づくもっともらしい現代神話を構成している．特に，のちに「ガベージコレクション」で主題として扱われる情報熱力学が色濃く出ていることに注目したい．

            物理学者にとって，熱力学は神話に近いかもしれない．熱力学は全宇宙・全時間に対して常に成り立つ最強の法則であり，人智を超えているという感覚がある．

        \subsection{p.87, l.14-}

            \begin{quote}
                「私たちの一日の大半は，〜」
            \end{quote}

            観測事実を持ち合い，解釈を提案しあい，吟味するという営みは，科学の本質である．「パリンプセストあるいは重ね書きされた八つの物語」の末尾にも，同様の記述が見られる．

            ところで，ここでキャサリンたちが行っている行為は，いままさに私が円城塔作品に対して試みていることとほとんど同じである．ここではなぜキャサリンたちがやたら多く存在するのだろうか，これは円城塔が好きだというヴァーリイ「バービーはなぜ殺される」に類似していないか，類似しているとすればそれはどのような意図が背景にあるのだろうか，そもそも作者の意図を考慮するべきだろうか，など．

        \subsection{p.89, l.16-}

            \begin{quote}
                「鹿を爪のようなものとしたために引き裂かれたキャサリンがおり，〜」
            \end{quote}

            観測と解釈を通じて，対象を何らかの形に“確定”することによって．キャサリンたちは世界を平定しようとする．ここには，地図作成と交易を目的としたデジタルゲーム『THE ATLAS』シリーズの影響が見える．また，デジタルRPGの元祖『Wizardry』の影響\footnote{ただし，円城塔は『Wizardry』の内容や有名なネタは知っているものの，実際にプレイしたことはないという（2026年5月8日開催の円城塔・前島賢『土人形と動死体』刊行記念対談での発言より．「小学生がそんなシビアなゲームをやるわけないでしょ」とのこと）．『Wizardry』に関する描写は『土人形と動く死体』「パリンプセスト」「良い夜を持っている」「世界でもっとも深い迷宮」に見られる．}も同時に指摘される．

            また，“鹿を爪のようなものとしたために”というのは，しかつめらしい（鹿爪らしい）とかけたギャグか．ただし，なぜしかつめらしいと引き裂かれるのかについては不明．

        \subsection{p.89, l.12-14}\label{duck}

            \begin{quote}
                「それはキャサリンAの形を持ち，キャサリンAであるかのように動くだけのキャサリンAにしかならないことが知られている．」
            \end{quote}

            ダックゲーム．円城塔作品には，この手の話が頻出する．

        \subsection{p.90, l.11-12}

            \begin{quote}
                「私たちはボーズ粒子ではありえないから．」
            \end{quote}

            キャサリンAとキャサリンBのそれぞれを中心とする同じ半径の円を完全に重ね合わせられないことに対する言明．これは定義より直ちに従う．

            この世の物質を細かく分解していくと，あるところでそれ以上分割できないような微粒子に還元される．そのような微粒子を素粒子(elemental particles)といい，素粒子は物質を構成するフェルミ粒子（フェルミオン）と，物質間に働く力を媒介するボーズ粒子（ボゾン）の2種類に分類される．

            このうち，ボゾンは同じ座標\footnote{厳密には，座標ではなく状態．フェルミオンにおいてはパウリの排他律によって禁止される．}にいくらでも詰め込むことが可能であることが知られている．このような現象を，ボーズ\hyphen アインシュタイン凝縮という．一方，フェルミオンはボーズ\hyphen アインシュタイン凝縮を起こさない．つまり，同じ時刻に同じ座標を占めることができない\footnote{ここから，柄谷行人が「意識と自然」\cite{krtn}で引用した，夏目漱石による有名な断片「二個の者がsame spaceヲoccupyスル訳には行かぬ．甲が乙を追い払うか，乙が甲をはき除けるか二法あるのみぢや．甲でも乙でも構わぬ強い方が勝つのぢや．」\cite{soseki}を連想するのはやや飛躍の感が強い．}．

            キャサリンは物質から構成されるので，キャサリンはフェルミオンの集団である．したがって，キャサリンは同じ時刻に同じ座標を占めることはできない．キャサリンが\alpha 粒子\footnote{陽子と中性子それぞれ2個ずつから成る粒子．ヘリウムの原子核と同一．放射線の一種\alpha 線は高速で飛行する\alpha 粒子．陽子と中性子はフェルミオンだが，それらから成る\alpha 粒子はボゾン．}である可能性はあるが，キャサリンが\alpha 粒子であると確信出来る記述が作中で確認出来ないため，キャサリンはフェルミオンだと仮定してよい．

        \subsection{p.91, l.1-10}

            \begin{quote}
                「例えば象を冷蔵庫に入れる冗句は，〜」
            \end{quote}

            前半は有名なジョークをそっくりそのまま説明しているもの．常識的には無茶苦茶なジョークであり，ルールはあると言えばあるが無茶苦茶である．

            後半は，そんな無茶苦茶なルールが罷り通り，ルールに穴があればたちまちつけ込まれるという話をしている．自然法則が法律のように策定され，運用されているというような気配がある．あるいは，観測者効果を過大評価した生半可な量子論の理解を意識した記述か．また，ここでいう麒麟は，首が折られたという描写から，空想上の生物の方ではなく，実在するジラフの方の麒麟のことだろう．

            自然法則において可能なことは必ず実現されるものであり，実際にそれは自然界で観測されうるし，われわれはそれを利用することが出来る，という物理学者が暗黙に共有している自然観が背景にある．この自然観は正しい\footnote{正しくなければ，工学という体系が存在しえない．}．

        \subsection{p.91, l.11-12}

            \begin{quote}
                「霧の向こうからピンク色の象が現れて私の横を通り過ぎ，〜」
            \end{quote}

                ピンク色をした象は，飲酒や錯乱によって正気を失った人間の見る幻覚の代表例としてよく用いられる．ここでも，観測者が正気を失っていることの象徴として登場しているものとしていいだろう．

                また．スタニスワフ・レム『ソラリス』で語り手が自身の正気を確かめていることとの関連も指摘される．のちの「バナナ剝きには最適の日々」では語り手が発狂しているらしいことにも注意したい．

                円城塔作品は，語り手の正気を自ら検証することが多い．これはただ気狂いを出して面白がっているのではないと擁護したい．実験物理学者にとって，実験の様子を観測する観測装置が正常に動作していることを確認することは最も重要な仕事である．かつて，CERN\footnote{欧州原子核研究機構(Conseil Européen pour la Recherche Nucléaire)．欧州各国が出資しあって設立された国際研究機関で，スイスとフランスの国境を跨いで地下に建設された直径9kmの巨大な加速器(Large Hadron Collider, LHC)が有名．}において，巨大な実験装置のうち，ケーブルが数cmたるんでいたことが原因で，標準模型を大きく外れる想定外の実験結果が観測されたことがあった．

        \subsection{p.93, l.8-10}

            \begin{quote}
                「「国立国語学研究所，異世界からの攻撃を受ける」〜」
            \end{quote}

            キレキレのギャグ3連発．

        \subsection{p.94, l.7}

            \begin{quote}
                「時事「モスクワ攻囲最終段階へ．円周率，3.75へ上昇」」
            \end{quote}

            この前後はまたもやパワーワードの連発．それにしても，円周率は定数であり，変動することは（常識的には）ありえないだろう．しかしそれが起こってしまっているので，物理学者としては円周率は変動するものと認めざるをえない．数学者がどう考えるかについてはまったくもって不明である．

            ところで，ディラックは，晩年に執着した大数仮説という妄想的主張において，あらゆる物理定数は実は時間変動していると主張した．本作の情景は，このディラックの大数仮説の数学定数への拡張としても捉えられる．

            \begin{clm}{ディラックの大数仮説}{}
                すべての物理定数は，時間変動する．
            \end{clm}

        \subsection{p.95, l.4}

            \begin{quote}
                「例えば円周率が3.8度線あたりで拮抗している〜」
            \end{quote}

            朝鮮戦争の停戦ラインで，北朝鮮と韓国の事実上の国境線（軍事境界線）となっている北緯38度線のこと．初期円城塔には，このようなやや不穏なギャグもよく見られる．

        \subsection{p.95, l.5-6}

            \begin{quote}
                「要するに円周率は何らかの収束していく数列の極限と見ることも可能であり，」
            \end{quote}

            正しい．14世紀の時点で，円周率は以下のような級数表示が可能であることが知られていた．

            \begin{thm}{マーダヴァ\hyphen ライプニッツの公式}{}
                \begin{align*}
                    \frac{\pi}{4} = 1 - \frac{1}{3} + \frac{1}{5} - \frac{1}{7} + \cdots = \sum_{n=0}^{\infty} \frac{(-1)^n}{2n+1}
                \end{align*}
            \end{thm}

            また，“インドの魔術師”シュリニヴァーサ・ラマヌジャンは，以下の公式を証明することなく直接言明した\cite{hardy,ramanujan}．

            \begin{thm}{ラマヌジャンの円周率の公式}{}
                \begin{align*}
                    \frac{1}{\pi} = \frac{2 \sqrt{2}}{99^2} \sum^{\infty}_{n=0} \frac{(4n)!}{4^{4n}(n!)^4} \frac{1103+26390n}{99^{4n}} \\ \frac{4}{\pi} = \sum^{\infty}_{n=0} \frac{(-1)^n (4n)!}{4^{4n} (n!)^4} \frac{1123+21460n}{882^{2n+1}}
                \end{align*}
            \end{thm}

            ラマヌジャンの円周率の公式は非常に奇妙な形ではあるものの，現在では数学的な証明が与えられている．

            \subsection{p.95, l.6-7}

            \begin{quote}
                「どこかの果てで円周率と一致すればそれでよく，手の届く有限部分での出来事などは関係ない」
            \end{quote}

            円周率の近似の収束速度のことを指す．先掲のマーダヴァ\hyphen ライプニッツの公式は，$n \to \infty$で$\displaystyle \frac{\pi}{4}$に収束するが，収束の速度は非常に遅い．有限部分での収束が遅くても，最終的に円周率に収束していれば（数学的には）どうでもいい，という主張は正しい．

        \subsection{p.95, l.8-10}

            \begin{quote}
                「当然円周率なんていうものは，収束する数列というだけには留まらず，全分野全方位的に支えられている怪物なので，〜」
            \end{quote}

            正しい．円周率は正規分布にも登場するので，あらゆる自然現象の背後にあると言って差し支えないだろう．ここで，正規分布は以下のように定義される\cite{tokei}．

            \begin{dfn}{正規分布（ガウス分布）}{}
                平均を\mu，分散を$\sigma^2$とする．\textbf{正規分布}とは，以下の確率密度関数で与えられる確率分布である．
                \begin{align*}
                    f(x) = \frac{1}{\sqrt{2 \pi \sigma^2}} \exp \left( - \frac{(x - \mu)^2}{2 \sigma^2}\right)
                \end{align*}
                ただし，$\sqrt{2 \pi \sigma^2}$は規格化定数．
            \end{dfn}

            正規分布は測定誤差の分布である．したがって，観測に基づく科学はすべて正規分布の影響下にある，つまり円周率の影響下にある．作中のように円周率が変動するのであれば，観測に基づく科学のいたるところで修正が発生し，玉突きで観測に基づかない科学にもその影響と修正が及ぶことになる．

            また，何らかの概念について，それと関連のあるもの同士で“支え合っている”という図式は，「ムーンシャイン」と共通する．

        \subsection{p.97, l.3}

            \begin{quote}
                「「勿論．必要にして充分なことだと思うね．」」
            \end{quote}

            梧桐のセリフ．数理科学において，“必要にして充分”という主張（同値であるという主張）は非常に強い主張であり，人の面前でこの主張をすることは，お前は同値関係すらわからない馬鹿なのか？というニュアンスがこもる．つまり，梧桐と私の議論がヒートアップしていることを示す．

        \subsection{p.98, l.11}

            \begin{quote}
                「「帰結は明らかだ」」
            \end{quote}

            “私”のセリフ．このセリフは，自明な議論を相手が理解してくれないときに発する．つまり，自分の説明が悪いのではなく，お前の頭が悪すぎるのがいけないのだ，という強い非難とともに匙を投げている．

        \subsection{p.100, l.5-10}

            \begin{quote}
                「私たちと同じような知識を持ち，〜」
            \end{quote}

            円城塔作品に頻出する，演繹的な言及の現れ．

            直後にある，運動方程式云々という話は，複雑系の考え方の現れ．

        \subsection{p.101, l.6}

            \begin{quote}
                「私たちは一を七に，〜」
            \end{quote}

            本作の進行はキャサリン128を追跡する語り手の経路をなぞっていることが明かされる．

        \subsection{p.103, l.15-}

            \begin{quote}
                「そんなにむきにならなくとも，〜」
            \end{quote}

            ここからしばらく，両立不能な2つの事柄を同時に追い求めたい，という話が続く．これは排中律の否定であり，先に披露されたNNP領域なる意味不明な記述とも関連するか．

        \subsection{p.107, l.3-4}

            \begin{quote}
                「武田騎馬軍団の六文字を映像で置き換えるのに必要とされるリソースは考えるだに気が遠くなる．」
            \end{quote}

            その通りという話なのだが，のちに兵馬俑がモチーフの「文字渦」を書くことや，小説の優位性として物語の高速探索が可能な点を挙げること\cite{kakuyomu}を考えると，これは円城塔がその初期から一貫して同じ主張を繰り返していることの証拠であるとも考えられる．

        \subsection{p.107, l.15-17}

            \begin{quote}
                「「例えば，全宇宙の素粒子の数に匹敵する数の火星人の侵略とか．〜」
            \end{quote}

            先の話の変奏ではあるが，ここではフレーゲ\hyphen ラッセルから始まる言語哲学の議論に接近しつつある．

        \subsection{p.108, l.11-}

            \begin{quote}
                「既にして充分蓄えられるだけ蓄えられて手付かずにいる自然言語コーパスと，定理自動証明理論．そいつらを一緒くたにかき混ぜて〜」
            \end{quote}

            現在のLLMの作り方そのもの．数学に関しては，おそらく難易度と需要の小ささが原因であまり大きな発展は見られなかったが，直近になって急速な発展が見られた．ただし，定理自動証明理論（自動定理証明）や定理証明支援系を用いたものではなく，あくまでも通常のLLMの延長として得られた．

        \subsection{p.109, l.8}

            \begin{quote}
                「そいつらに内面があるのかどうかという深遠くさい問題」
            \end{quote}

            2026年現在，まさに問題になっているもの．Claudeを擁するAnthropicのCEO，ダリオ・アモデイは，LLMに意識がないとは言い切れないと主張している\cite{amodei1}．これはAnthropic社内で議論されていた内容\cite{amodei2,amodei3,constitution}の吐露と見られる．一方で，テッド・チャンは「LLMは応用統計学と呼ぶべきものであり，LLMに意識があるなどということはない」と繰り返し主張している\cite{chiang1, chiang2}．とはいえ，チャンはやや不合理なまでの自由意志の信奉者であり，冷静な評価とは言い難いところも指摘される．

            円城塔作品では，\ref{duck}節でも見た通り，ダックゲームが頻繁に描写され，またウルフラムの提唱がやや濫用されている．また，円城塔は『ムーンシャイン』のあとがきにおいて，次のように述べている\footnote{ここで述べられている着想は，のちに『土人形と動死体』に結実する．}\cite{moonshine}．

            \begin{quotation}
                人間の歴史の大部分は「自分は人間である」と主張する相手を「お前は人間ではない」と否定することで続いてきた．(中略)機械が実際，人間としか思えない段階に達した場合に想像されるのは「人間扱いされるようになった機械」と「機械扱いされるようになった人間」の併存であり，それはすでに様々な形で生まれていると見ることもできる．
            \end{quotation}

            これらのことから，円城塔は，意識と判定する基準が甘い傾向\footnote{むしろ，何者かによって作られてしまったもの（物語の登場人物，物語，LLM）を世界内存在として捉え，存在することの苦悩に同情する傾向が指摘される．例として「墓の書」「ローラのオリジナル」が挙げられる．}にあることが推測される．つまり，円城塔はLLMに意識があると考えることに抵抗は少ないものと考えられる．

        \subsection{p.110, l.5}

            \begin{quote}
                「「お前にもそうしてきた記憶があるだろう」」
            \end{quote}

            キャサリンたちがその誕生直後から過酷な“世界”に投げ込まれてあるのに対して，“私”は親の庇護下で積み木や粘土などの適度な刺激を与えられて成長してきただろう，という問いかけ．これは発達の過程を複雑系として捉えようという研究を背景としているか．ここには，ハイデガーの“世界内存在”の影響も指摘される．

            また，複雑系の研究者は，子供の発達を観察することを好む．円城塔作品では『プロローグ』にその傾向が見える．

        \subsection{p.112, l.17}

            \begin{quote}
                「これは私たちに与えられた言語コーパスの癖」
            \end{quote}

            キャサリンの1人の独白．LLMの構成方法とその解釈として正しい．LLMは，宇宙エレベーターに関する文章を学習することで，“宇宙エレベーター”という単語の前後に配置されるべき単語の確率分布を獲得する．

        \subsection{p.113, l.10-17}

            \begin{quote}
                「この円柱や〜」
            \end{quote}

            そのまま読めばいい．

        \subsection{p.115, l.10-11}

            \begin{quote}
                「今の数学が成長する以前に存在し，今の数学に圧倒されて滅亡の淵に追い込まれた萌芽的論理構造．」
            \end{quote}

            作中では，このような現在の数学とは異なる数学のことを代替数学と呼んでいる．確かに，数学にも人間の認知機構，さらに言えば人間の脳神経構造に由来する偏りがあるかもしれない．しかしながら，通常の数学ではあまり意識されないが，数理論理学では数学が成り立つための条件の探索が進んでおり，このような心配は払拭されている．

            数学のような形式科学が（その発展の途中では）人間の認知機構に由来する偏りがあると考えるならば，社会構造に由来する偏りについても検討されなければならないだろう（社会構築主義）．この社会構築主義的な姿勢は，のちに「ニュー・ナンバーズ」で明確に示される\footnote{ソ連におけるルイセンコ学説の事例は，政治的正しさが自然科学に影響を与えた社会構築的な事例として教訓にすべきである．しかしながら，十分発展した数学・物理学において，社会構築的な要素が残るとは考えにくい．多くの物理学者は，数理科学が社会構築的なものであるとは考えていない．物理学者の大野克嗣は．村上陽一郎に代表される，日本の科学論者がしばしば展開する粗雑な社会構築主義的な議論を批判している\cite{oono}p.194．生物学者からも，生物学に対して科学論者が寄せる粗雑な社会構築主義的な主張に対する批判がある\cite{hasegawa,minaka}，}．

            また，代替数学の探索は，可能なすべての数学の探索の一環と言える．この立場からすれば，このような探索は，可能なすべての経路の探索を試みる解析力学や，可能なすべての生命の探究を試みる理論生物学，可能なすべての文字列の考察に取り組む円城塔の姿勢と一貫している．

        \subsection{p.115, l.14-}

            \begin{quote}
                「私たちは岩壁から仏を切り出している．〜」
            \end{quote}

            「ガベージコレクション」っぽい．直接的にはバーミヤンの石仏を指すか．

        \subsection{p.117, l.2-5}

            \begin{quote}
                「わたしの名はキャサリンA．〜」
            \end{quote}

            「Self-Reference ENGINE」のクライマックス，「私の名はSelf-Reference ENGINE，〜」が直ちに連想される．

        \subsection{p.117, l.6}

            \begin{quote}
                「こんなこともあろうかと」
            \end{quote}

            『宇宙戦艦ヤマト』に登場する技術者，真田志郎の有名な台詞とされるもの．実は一言一句合致する台詞は存在しないことが知られている．ここでは，どんな不測の事態に対してもなぜか用意周到であるというお約束をギャグとして挿入しているか．

        \subsection{p.118, l.1}

            \begin{quote}
                「\ruby{夜久毛多都}{やくもたつ}〜」
            \end{quote}

                記紀において，スサノオノミコトが詠んだとされる歌．紀貫之『古今和歌集』仮名序によれば，日本最古の和歌とされる\footnote{「人の世となりて，すさのをの命よりぞ，三十文字あまり一文字はよみける．」\cite{kokin}p.10}\cite{kokin}． 現代神話としてのSFをより強調する．

                直後にあるのは万葉仮名で書かれたスサノオの歌を直訳したもの．「〈ゲンジ物語〉の作者，〈マツダイラ・サダノブ〉」やラフカディオ・ハーン『怪談』の翻訳，ひいては『去年，本能寺で』につながる．

        \subsection{p.118, l.5-8}

            \begin{quote}
                「多分，どこかの方向に気の違った流行みたいなものだったのだ．〜」
            \end{quote}

            具体例としてタートルネック，ルーズソックス，長裃の袴，サーベルタイガーの牙などが挙げられていることから，時間発展のうちに一部の器官が本来求められる性能を超えて長大化した現象を指している．

            サーベルタイガーを例にとる．サーベルタイガーの牙は巨大化しすぎたため，強い力で獲物を噛むと牙が折れてしまう．これを裏付けるように，サーベルタイガーの代表的な種，スミロドンの咬合力は同じ体格である現代のライオンの$\displaystyle\frac{1}{3}$程度しかなかったと推定されている\cite{sabercat}．これはサーベルタイガーの（立派すぎる）牙に対して一般的に抱かれるイメージとは異なる．

            また，いわゆるサーベルタイガーはネコ科で巨大な犬歯を持つ種を指すが，ネコ科以外の種にも同様に巨大な犬歯を持つサーベルタイガーに近い特徴を持つ種が存在していたことが知られている\footnote{有袋類（カンガルーなど）に近いティラコスミルスなど．}．つまり，犬歯の巨大化は進化史上一度だけ発生した例外的な事例ではなく，複数回発生した普遍性を持つ現象であることが示唆される．ある時代に生息していたサーベルタイガーが絶滅したり，そもそもその時代のその地域にサーベルタイガーが存在しなかったりすると，そのニッチ\footnote{環境的地位のこと．ある環境において，ある種が占有する資源・生息地などを指す．}を埋めるように新たなサーベルタイガー的な種が出現した．この解釈として，巨大な牙は一定程度生存に有利に働くが，ある度合いを超えて巨大化しすぎると不利に働いて絶滅に向かう，というものが考えられる．このような，一時期は有利だが，その先には絶滅しかないような危険な進化論的アトラクタのことを，円城塔は“気の違った流行”と表現しているものと推察される．実際，ルーズソックス\footnote{2025年においてはむしろ復活していたことに注意．やはり空いたニッチを埋めるように進化するか．}，長裃，サーベルタイガーは現在では絶滅している．

            この解釈に対して，定向進化説による異説を提出することも出来る．つまり，ネコに似た種が進化の袋小路に入った場合，犬歯が巨大化するという末期的な容態を呈し，やがて絶滅する．

            このような定向進化説的な解釈は，実際に異常巻きアンモナイトと呼ばれる一部のアンモナイトの解釈として用いられた．異常巻きアンモナイトは，いわゆる普通のアンモナイトに比べて巻きが緩く，3次元的な広がりをもつという特徴があり，アンモナイトの中では末期に出現した種である．“異常巻き”と表現される通り，緩やかな巻きが種としての活力が衰退していることを連想させたため\footnote{異常巻きアンモナイトの画像を適宜参照いただけると，“末期的な容態”という表現にも納得いただけるだろう．}か，かつては進化の袋小路に入った末期的な容態であると考えられていたが，現在では否定されている．アンモナイトの巻き方は数学的に記述することが出来，その巻き方を特徴付ける変数を操作することで，異常巻きアンモナイトも普通のアンモナイトも全て統一的に記述される\cite{okamoto,matsumoto}．この異常巻きアンモナイトの数理モデルは，生物の容態が数学的に規定されている事例として，イモガイの貝殻の模様がセルオートマトンで記述出来ること，タテジマキンチャクダイの体表の模様がチューリングパターンで記述出来ることと並んで有名である．なお，異常巻きアンモナイトの化石は日本で多く産出し，研究が盛んである．

        \subsection{p.120, l.15-17}

            \begin{quote}
                「一つ．ラマルクの定向進化説は間違っていたということは間違っていた．二つ．ルイセンコはどう裏返してみても嘘吐きだった．三つ．ケストラーは第十三支族について言及した後，ホテルの一室で変死した．」
            \end{quote}

            ラマルク，ルイセンコ，ケストラーはいずれもダーウィン\hyphen メンデルの進化論（より厳密には総合進化説）と相反する主張を行った人物．ダーウィン本人の主張にはのちに否定された主張も多く\footnote{例えば獲得形質遺伝説．}，すべてを紹介すると混乱が生じるので，以下の重要な点3点に注目する\cite{sarashina}．

            \begin{clm}{ダーウィンの進化論の要点}{}
                \begin{enumerate}
                    \item 進化説\\生物は進化しうる．
                    \item 自然選択説\\進化は自然選択によって起こる．
                    \item 分岐説\\生物は1つの種から複数の種に進化しうる．
                \end{enumerate}
            \end{clm}

            ジャン-バティスト・ラマルクは，チャールズ・ダーウィンより半世紀ほど前の18世紀から19世紀に活動した博物学者\footnote{当時は生物学がまだ独立していなかった．そもそも，生物学(biology)という言葉自体がラマルクによるものである．}．ラマルクは1800年ごろに生物は進化するという主張（進化論）を認めていることを告白し\footnote{進化論を唱えたのはラマルクが最初ではない．ラマルクが『動物哲学』を通じて普及させたのは確かなことなのだが，ラマルクのほか，チャールズ・ダーウィンの祖父エラズマス・ダーウィンやゲーテも進化論に近い考えを表明しており，やや遡ってカントやディドロ，さらには古代ギリシアのアナクシマンドロス，エンペドクレス，アリストテレスも種は変わりうるものであるとの認識を持っていた\cite{sarashina}．}，1809年に出版した啓蒙書『動物哲学』\cite{lamarck}では当時流通していた進化論をまとめ，大衆向けに概説した．

            なお，本作ではラマルクが定向進化説を主張していたということになっているが，これは誤りである．ラマルクの主張を整理すると，以下のようになる\cite{mori, kawada}．

            \begin{clm}{ラマルクの進化論の要点}{}
                \begin{enumerate}
                    \item 前進的発達説\\生物は内在する発達傾向によって，単純なものから複雑なものに進化する．
                    \item 自然発生説\\単純な生物は自然に発生する．
                    \item 用不用説\\生物は獲得形質の遺伝によって，環境変化に適応するように進化する．
                \end{enumerate}
            \end{clm}

            ラマルクの主張の中核は，前進的発達説である．これは進化の根源的要因を生物の内在的な発達傾向に求め，生物は自ら単純なものから複雑なものに進化するという考えである．一方で，ラマルクは種は進化によってより複雑なものに進化するとしつつも，ある1つの種から複数の種に分岐するような進化を考えなかった．むしろ，単純な生物は自然に発生し，複雑な生物に向かって直進的に進化すると考えた（自然発生説）．つまり，複雑な種（哺乳類）は遠い過去に発生して進化した種であり，単純な種（昆虫）は最近発生した種であるとした．また，生物は自らの意志によって自身の形質を改良し，その形質を子孫に伝えることで，より環境に適応するように進化する（用不用説・獲得形質遺伝説）\footnote{社会ダーウィン主義（社会ダーウィニズム）として知られる思想は，実はダーウィンよりもラマルクに近い．これには，ダーウィンとラマルクの両者から影響を受けた哲学者・社会学者，ハーバート・スペンサーの影響が非常に大きい．スペンサーは”適者生存”という言葉を生み出し，適者生存を社会に対して適用した．}．

            定向進化説（系統発生説とも）は，1890年代に新ラマルク主義（ネオラマルキズム）の古生物学者，テオドール・アイマーらを中心に提唱された説である．その主張は，進化とはその進化によって得られた形質が適応的かどうかに関わらず，一旦進化が進んだ方向に向かって種全体が一定の方向に進化する，というものである\cite{kawada,shikama}．その具体例として，進化に伴って指の本数が減少し，体躯が大きくなるようになったウマがしばしば用いられる．定向進化説は，その種に内在する進化の方向性を仮定する点でラマルクの前進的発達説・用不用説に似ているが，ラマルクが進化の方向性は適応に有利な方であると考えたのに対して，定向進化説は適応の有利不利を問わない\footnote{定向進化説が適応の有利不利を問わないのは，おそらくこの説が主に古生物学者によって形成・支持されたことによる．定向進化説は，ある傾向に従って進化しているように見える化石が存在することを根拠の1つとするが，このような事例は自然選択説など他の説でも十分に解釈出来，説得力に欠ける．定向進化説は分子生物学における実証に耐えないという評価がなされ，現在では完全に放棄されているが，古生物学では長く有力な説であった．教科書の1章を割いて定向進化説を詳説する\cite{shikama}はその好例である．}．

            また，定向進化説（系統発生説）に関連して，反復説について触れたい．反復説とは，個体発生（個体が受精卵から胎児に変化していく過程）は系統発生（その個体が属する種の進化の過程）を同じ順番で繰り返すというもので，19世紀の生物学者・哲学者エルンスト・ヘッケルによって提唱された．反復説は現在では否定されているが，その議論の過程で多数の有益な知見が得られたことや，社会科学に多大な影響を与えたこと，また同時代の文学や今日に至るまでの教育カリキュラムにも影響を与えたことが知られている．

            例えば，スティーヴン・J・グールドは，『個体発生と系統発生』において，J・G・バラード『沈んだ世界』に反復説の影響が見えることを指摘している\cite{ontogeny}p.181．%佐藤恵子『ヘッケルと進化の夢』
            また，「Boy's Surface」「Your Heads Only」「Gernsback Intersection」に引用されている宮沢賢治\footnote{このほか，「ドルトンの印象法則」において『春と修羅』序文が実験操作の対象として用いられている．}にも反復説の影響が指摘されている\footnote{青木（秋枝）美保は，「青森挽歌」にある“ヘッケル博士”の解釈について，先行研究を整理分類している．青木によれば，先行研究による解釈は，宮沢賢治がヘッケルの主張した霊魂死滅説に同意している，あるいは反対に霊魂不滅説に同意しているという2通りに分類されるという\cite{akieda}．これに対して，浜垣誠司は，霊魂死滅説・霊魂不滅説ではなく，同じくヘッケルが主張した反復説への同意であると解釈する\cite{hamagaki}．}．

            円城塔は，自身がSF考証・脚本を担当した『ゴジラS.P』におけるゴジラについて，個体発生が系統発生する生物というコンセプトとしたとしており，反復説は円城塔研究においても検討されるべきである．

            %ダーウィン前夜の進化論

            トロフィム・ルイセンコは，スターリン時代に活躍したソ連の農学者・生物学者・農業指導者．ルイセンコは農民の家庭に生まれ，キエフのポルタワ園芸専門学校を卒業し，ギャンジャ\footnote{現在のアゼルバイジャン西部の都市．東部のカスピ海沿岸の首都バクーに次ぐ第2の都市．}の育種試験場に就職した\cite{nakamura,medvedev}．ギャンジャでは，1926年から翌年にかけて，綿の作付けに先立ってエンドウマメを冬季に栽培する研究に従事し．独創的ではないかもしれないが，冬でも暖かいアゼルバイジャンの気候に合った有益な農法を提案して名声を得た\footnote{ダーウィン\hyphen メンデル以来の正統遺伝学に基づき多大な功績を残したソ連の生物学者で，のちに反ルイセンコ派の筆頭となるニコライ・ヴァヴィロフは，この研究成果を知ってルイセンコに興味をもち，自身の所属する全ソ連植物栽培研究所に招聘しようとしたが，所内から反対意見が挙がり実現されなかった．なお，ヴァヴィロフが部下に命じてルイセンコについて調べさせた報告書では，「ルイセンコは教育がなく，自尊心の強い人間であるが，有能な実験家であることに間違いない」とある\cite{medvedev}．}\cite{medvedev}．ルイセンコが一躍有名になったのは，1930年代に取り組んだ，小麦の春化処理（ヤロビザーツィヤ）に関する研究だった\footnote{なお，春化処理は1929年に初めて論文として出版されているが，その内容と実験手法は粗末であり，奇しくも同じ会議でニコライ・マクシモフが出版した，より学術的に洗練され，しかも先行研究も上手くまとめられた論文の陰に隠れて目立たなかった．そもそも，春化という現象自体は1857年にアメリカの農業指導者，ジョン・ハンコック・クリッパートによって報告されている\cite{medvedev}．}．

            植物には，品種によって種子の発芽に必要な温度が異なる性質があり，このような性質を播性という．小麦には大きく分けて春播小麦と秋播小麦の2種類があり，例えば秋播小麦を春に播いても成長しないのだが，生長過程で適当な温度変化を加えると春でも成長し，収穫を得られるようになる．これを春化処理という．ルイセンコは，春化処理が獲得形質の遺伝によるものであると独自に解釈し，ダーウィン\hyphen メンデルの正統遺伝学を否定した．

            ルイセンコは春化処理によって小麦の収量が増大すると主張し，また自説をマルクス\hyphen レーニン主義と合致するように彩った．一方，自説と相反する正統遺伝学とそれに基づく従来の農学はブルジョワ科学であるとして批判し，スターリンによる後ろ盾を得て，反対者を政治的・学術的に，あるいは実際に抹殺した\footnote{前述のヴァヴィロフは，1940年に逮捕され，1943年にシベリアで獄死した．この他，遺伝学者カルペチェンコ，細胞学者レヴィツキー，植物学者ゴヴォロフら正統遺伝学に基づくソ連の有能な生物学者が多数獄死させられ\cite{medvedev}，ソ連の正統生物学は壊滅した．}．

            ルイセンコ学説に基づく農業（ミチューリン農法\footnote{ソ連の有能な育種家，イヴァン・ミチューリンに因む．ミチューリンは正規の高等教育を受けていないが，園芸学を独学し，寒冷なロシアの気候に合った果樹品種を多数開発した．スポークスマンのイサーク・プレゼントと共に，ルイセンコは自身をミチューリンの後継者であると主張し，またソ連共産党がルイセンコ学説をミチューリン\hyphen ルイセンコ遺伝学と呼ぶことを認めたため，ルイセンコ学説に基づく農業をミチューリン農法と呼ぶようになった．}）がソ連で広く実践された\footnote{ミチューリン農法は東側諸国に積極的に輸出され，各国に被害をもたらした．特に中国では大躍進政策に，北朝鮮では主体農法に組み込まれ，甚大な被害をもたらした．西側諸国においては英国のJ・D・バナール以外には当然無視されたが，日本ではヤロビ農法と称され左翼知識人と農家を中心に紹介・実践が行われた．日本における状況は\cite{nakamura}を参照されたい．}が，当然ながら作物の収量は悪化し，ソ連の食料供給にダメージを与えた．このように収量悪化の報告が多数届けられたこと，最大の庇護者スターリンが死亡したことで，1950年代半ばにルイセンコ学説の求心力は衰退した．また，ヴァヴィロフの死後に反ルイセンコ派の筆頭となった遺伝学者，ニコライ・ドゥビーニンらを中心として，イーゴリ・タム\footnote{理論物理学者．専門は原子核理論．1953年にチェレンコフ効果の理論的解釈によってノーベル物理学賞受賞．}，ピョートル・カピッツァ\footnote{実験物理学者．専門は低温物理学．1978年にヘリウム4の超流動状態の発見などの功績によってノーベル物理学賞受賞．}，レフ・ランダウ\footnote{理論物理学者．専門は物性理論．1962年にヘリウム4の超流動状態の理論など凝縮系における功績によってノーベル物理学賞受賞．物理学の有名な教科書，ランダウ\hyphen リフシッツのランダウ．}，アンドレイ・サハロフ\footnote{理論物理学者．専門は宇宙物理学，原子核理論．1975年にノーベル平和賞受賞．40-50年代のソ連における原爆・水爆開発の中心人物の1人だが．核実験による環境への悪影響が明らかになると転向し，人権擁護活動に尽力した．}，ヴィタリー・ギンツブルク\footnote{理論物理学者．専門は物性理論．2003年に超伝導と超流動に関する理論への貢献によってノーベル物理学賞受賞．超伝導のマクロ理論であるギンツブルク\hyphen ランダウ理論はランダウとの共同研究の成果．}，ムスティスラフ・ケルディシュ\footnote{応用数学者・計算機科学者．ソ連における計算機開発・航空工学・ロケット工学・核開発の重要人物．}，ミハイル・ラヴレンチェフ\footnote{数学者・応用数学者．ケルディシュの指導教官．ケルディシュとの共同研究のほか，爆発過程の流体力学的解析でも功績がある．}，セルゲイ・ソボレフ\footnote{数学者・計算機科学者．ソボレフ空間，ソボレフ不等式で有名．}，アレクセイ・リアプノフ\footnote{数学者・応用数学者・計算機科学者．リヤプノフ，リャプノフ，リャプーノフとも．ソ連におけるサイバネティックスの創始者．リアプノフ関数，リアプノフ指数など力学系で多大な功績のある数学者，アレクサンドル・リアプノフの甥．生物学を志す2人の娘のためにドゥビーニンを招いて私的なゼミを行なっていたところ，ソ連の生物学界がルイセンコ学説に汚染されていることに気づき，反ルイセンコ運動に参画した．}を含む297名によるルイセンコを糾弾する書簡「300人の手紙」が1955年10月と1956年2月\footnote{フルシチョフがいわゆるスターリン批判を行なったのも，この2月のことである．}の2度にわたってソ連共産党中央委員会幹部会宛に発送されたが，いずれも黙殺された\cite{ichikawa1,ichikawa2}．フルシチョフはルイセンコ学説を復興させたが，西側諸国で分子生物学の知見が急速に蓄積されていったことでルイセンコ学説の支持者は減少していき．1964年にフルシチョフが失脚するとともにルイセンコ学説はソ連の公式見解から除かれた．

            アーサー・ケストラーは，ハンガリー出身のユダヤ系作家．ウィーン工科大学を卒業後ジャーナリストとなり，自身や知人の経験を元にしたスターリン体制に批判的な小説『真昼の暗黒』で名声を得た．自然科学の知見を背景に，科学史や科学思想に関する著作も多いが，陰謀論や疑似科学\footnote{ケストラーは，心理学や生物学に関する還元主義（心理学は脳などの生物学的な器官に還元され，生物学は化学や物理学へと還元されるという主張）を批判し，否定しようとするあまりしばしば疑似科学的な主張に手を染めた．『機械の中の幽霊』や『ホロン革命』がその代表例で，ケストラーの著作はニューサイエンスやニューエイジ思想に大きな影響を与えた．また，『機械の中の幽霊』は士郎正宗『攻殻機動隊』『アップルシード』に影響を与えている．ここで，円城塔が士郎正宗『攻殻機動隊』の忠実な再現を謳うモコちゃん『攻殻機動隊』でシリーズ構成・脚本を担当していることに注意したい．}への傾倒が強く見られる．

            特に，『サンバガエルの謎』\cite{koestler1}はその陰謀論的傾向が強く見える．同書は1920年代に獲得形質遺伝説を実験的に証明したと主張した生物学者，パウル・カンメラーに関するノンフィクションである．カンメラーはサンバガエルを用いて獲得形質遺伝説を裏付ける結果を得たが，生物学者の多くから批判され，さらに批判派の生物学者による検証で実験不正を指摘され，最終的に自殺した．ケストラーはカンメラーと獲得形質遺伝説を熱心に擁護し，正統生物学を過度に信奉する者たちによって不当に貶められた被害者としてカンメラーを描いた．

            また，既に紹介した通り，ケストラーはユダヤ人であり，ユダヤ人に関する著作も発表している．1976年に出版された『ユダヤ人とは誰か』\cite{koestler2}において，ケストラーは，アシュケナジム\footnote{ディアスポラののち，ドイツ語圏あるいは東欧に定住したユダヤ人の子孫．}はユダヤ教に改宗したテュルク系集団であり，民族的なユダヤ人ではないと主張した（アシュケナジムのハザール起源説）．同書ではアシュケナジムを“第十三支族”と呼び，現在のユダヤ人には13番目の支族が存在するとして，ユダヤ人に対する人種差別の根拠を否定しようと試みた．同書はハザール起源説を大衆に広めるきっかけとなったが，遺伝学・言語学・歴史学による検証によって，現在ではハザール起源説は科学的に否定されている．

            最終的に，1983年にケストラーはホテルの一室で妻と共に遺体となって発見された．服毒自殺であると見られている．この事実が，本作でのケストラーに関する描写の根拠である．

            さて，本作の描写の解釈について．ラマルクについては，直前で描かれた，無闇に長大化した算盤の件を定向進化説に絡めて解釈したものである．ルイセンコは疑似科学に対する嘲笑であり，ケストラーはやや飛躍の強い連想の現れである．

            進化論が本作に登場すること自体については，題名の元ネタであるバッハ『ゴルトベルク変奏曲』を有名にしたピアニスト，グレン・グールドから，進化論に関する著作が多くある生物学者，スティーヴン・ジェイ・グールドを連想したのではないかと推察される．あるいは，進化論のグールドからピアニストのグールドに連想した可能性も指摘される．

        \subsection{p.121, l.12-13}

            \begin{quote}
                「人間は額に命令文を記述してやることで，巨人を操ることが可能である．」
            \end{quote}

            ゴーレムのこと．SFではテッド・チャン「七十二文字」で扱われたほか，のちに『土人形と動死体』で中心的な題材となる．

        \subsection{p.121, l.14-15}

            \begin{quote}
                「電子回路で構成された百手巨人」
            \end{quote}

            CPU(Central Processing Unit)のことか．

        \subsection{p.121, l.16-17}

            \begin{quote}
                「計算できるものは全て計算できるのだ．〜」
            \end{quote}

            チャーチ\hyphen チューリングの提唱に至るまでの議論を指す．チャーチ\hyphen チューリングの提唱はその名の通り人間が提唱した概念であり，数学的に証明出来ない概念である．

            1900年に，数学者ダーフィト・ヒルベルトは，パリで開催された国際会議において，「数学の問題」と題した講演を行なった．講演において，ヒルベルトは最も重要な数学上の未解決問題を10個紹介し，のちに23個に増補したリストを著書に掲載した．この“ヒルベルトの23の問題”は，当時最も偉大な数学者がわざわざ示した未解決問題のリストということで大きな話題になった\footnote{実際に非常に大きな成果に繋がった問題も多い．最も有名なところでは，第1問題は連続体仮説のZFC公理系からの独立性の証明，第2問題は第一不完全性定理・第二不完全性定理に繋がり，第6問題は物理学の公理化（未解決），第8問題はリーマン予想（未解決），という間違いなく有益だが至上の難問を扱っている．}\cite{hilbert}．その中の第10問題は“ディオファントス方程式が整数解をもつかどうかを判定するアルゴリズムは存在するか”というものであったのだが，当時はアルゴリズムという概念がまだ曖昧であり，まずはアルゴリズムを数学的に定式化する作業が行われた\cite{sipser2}．なお，ディオファントス方程式は以下のように定義される．

            \begin{dfn}{ディオファントス方程式}{}
                \textbf{ディオファントス方程式}とは，整数解または有理数解のみをもつような方程式であって，以下のように表記される整係数多変数不定方程式である．
                \begin{align*}
                    \Sigma a_{i_1,...,i_n} x_1^{i_1}...x_n^{i_n} = 0, \, a_{i_1,...,i_n} \in \mathbb{Z}
                \end{align*}
            \end{dfn}

            ピタゴラス方程式$x^2 + y^2 = z^2$はディオファントス方程式の代表例である．

            1936年に，アロンゾ・チャーチとアラン・チューリングがそれぞれ独立にアルゴリズムを定義した．ここでチャーチはラムダ計算\footnote{現在の多くのプログラミング言語に実装されている，無名関数（ラムダ式）の数学的基盤．}という記法を導入し，チューリングはチューリング機械を導入した．これらは一見して別物だったが，これらの定義は等価であることが示され，またラムダ計算・チューリング機械以外にも等価な計算モデルが複数発見された．これを受けて，これらの計算モデルはすべて等価であり，これらの計算モデルが実行可能なものを計算可能であるとみなすことになった．この取り決めのことを，チャーチ\hyphen チューリングの提唱という．

            \begin{ths}{チャーチ\hyphen チューリングの提唱}{}
                その計算が\textbf{計算可能}であるとは，その計算がチューリング機械によって構成可能なことである．
            \end{ths}

            ヒルベルトの第10問題は，当時22歳のユーリ・マチャセヴィッチによって1970年に否定に解決された．また，この発見の副産物として，すべての素数を出力しうる多項式（素数表現多項式）が得られた．これは当初24変数37次多項式として得られ，10変数11281次という変種も報告されている．これら素数表現多項式のなかで最も有名なものは，以下のジョーンズ\hyphen 佐藤\hyphen 和田\hyphen ウィーンズ多項式である\cite{seisuutan1}．

            \begin{dfn}{ジョーンズ\hyphen 佐藤\hyphen 和田\hyphen ウィーンズ多項式}{}
                以下の多項式の変数に非負整数を代入して得られる値全体の集合と，正の整数全体の集合の共通部分は，素数全体の集合に一致する，つまり，以下の多項式に非負整数を代入したとき，得られた値が正の整数であるならば，それは必ず素数である．
                \begin{align*}
                    (k &+ 2)(1 - (wz + h + j -q)^2 - ((gk + 2g + k +1)(h + j) + h - z)^2 \\ &- (2n + p + q + z - e)^2 - (16 (k + 1)^3(k + 2)(n + 1)^2 + 1 - f^2)^2 \\ &- (e^3 (e +2)(a +1)^2 + 1 - o^2)^2 - ((a^2 - 1)y^2 + 1 - x^2)^2 \\ &- (16 r^2 y^4 (a^2 - 1) + 1 - u^2)^2 \\ &- (((a + u^2(u^2 - a))^2 - 1)(n + 4dy)^2 + 1 - (x + cu)^2)^2 \\ &- (n + l + v - y)^2 - ((a^2 - 1) l^2 + 1 - m^2)^2 - (ai +k + 1 - l - i)^2 \\ &- (p + l (a - n - 1) + b(2an + 2a - n^2 -2n -2) - m)^2 \\ &- (q + y(a - p - 1) + s(2ap + 2a - p^2 - 2p -2) -x)^2 \\ &- (z + pl(a - p) + t(2ap - p^2 -1) - pm)^2)
                \end{align*}
            \end{dfn}

            ジョーンズ\hyphen 佐藤\hyphen 和田\hyphen ウィーンズ多項式は26変数25次多項式であり，変数としてアルファベットをちょうど使い切ることで知られている．

            また，素数表現多項式とも異なる別の副産物として，ディオファントス方程式と数学基礎論の深淵な関係が知られている．ヒルベルトの第10問題の解決に大きな貢献のあったマーティン・デイヴィス，ジュリア・ロビンソン，ユーリ・マチャセヴィッチ\footnote{この3人のほか，デイヴィスとともにゲーデル数化など論理学方面から第10問題の検討に大きく貢献した人物として，哲学者・論理学者ヒラリー・パトナムが挙げられる．しかしながら，ほどなくパトナムの関心は論理学・分析哲学の話題からかつては“擬似問題”と見做していたものに移ってしまい，以降数学・論理学に関する業績は途絶えた．なお，円城塔が好んでいる数学者レイモンド・スマリヤンの指導教官はパトナムである\cite{hilbert}p.140．}の大部の共著論文が存在しており，そこではある特定のディオファントス方程式\footnote{\cite{legends}で具体的に与えられている．数学上の問題をゲーデル数化することで，ディオファントス方程式に還元させて得られる．}が解を持たないとき，かつそのときに限り，肯定的に解決されるような問題が存在し，しかもそのような問題の中に，フェルマーの最終予想，ゴルトバッハ予想，ZFCにおける連続体仮説の証明可能性，リーマン予想などが含まれていることを示した\cite{hilbert,legends}．つまり，ヒルベルトの第10問題がヒルベルトの意図通り肯定的に解決されていたならば，これらの難問はゲーデル数化を通じてディオファントス方程式に帰着させることで，その解の存在を確かめることで証明・反証が可能となっていた．

        \subsection{p.122, l.6-}

            \begin{quote}
                「実際のところ，自然を支配する方程式が知られたとして，それを実際に計算することができるとは限らない．〜」
            \end{quote}

            この主張は正しい．ある現象が数学的に記述できたとしても，それが数学的に解くことができないということも十分考えられる．先ほどのディオファントス方程式には解が存在しない場合もある．また別の例として，流体力学が挙げられる．古典力学における運動方程式（ニュートン方程式）$\bm{F}=m\bm{a}$に対応する支配法則として，流体力学ではナヴィエ\hyphen ストークス方程式が知られているが，その解が物理学的・数学的に“いい振る舞い”をするかどうかは未解決となっており，これはミレニアム懸賞問題にもなっている．

            \begin{dfn}{ナヴィエ\hyphen ストークス方程式}{}
                ニュートン流体，つまり非圧縮性をもち，等質であるような流体に関するナヴィエ\hyphen ストークス方程式は以下の通り．ただし，$\rho$は流体の密度，$\bm{v}(\bm{x}, t)$は流体の速度，$\bm{p}(\bm{x}, t)$は圧力，$\nu$は粘性係数，$\bm{f}$は流体にかかる単位体積あたりの外力．
                \begin{align*}
                    \rho \bigg\{ \frac{\partial \bm{v}}{\partial t} + ( \bm{v} \cdot \nabla) \bm{v} \bigg\} = - \nabla \bm{p} + \nu \nabla^2 \bm{v} + \rho \bm{f}
                \end{align*}
            \end{dfn}

            \begin{prb}{ナヴィエ\hyphen ストークス方程式の解の存在と滑らかさ}{}
                $3+1$次元時空において，初期速度を与えたとき，ナヴィエ\hyphen ストークス方程式の解となる速度ベクトル場と圧力のスカラー場が存在して，しかも双方とも滑らかで大域的に定義されるか．
            \end{prb}

            また，“正確に解く”という正確さの精度にも依存する．摩擦や空気抵抗を無視した理想的な系では慣性の法則に従って等速直線運動を行う物体を正確に記述することもできるだろうが，現実的な系では非常に困難である．

        \subsection{p.122, l.6-17}\label{deathofgodel}

            \begin{quote}
                「大半の方程式は解を持つ．〜」
            \end{quote}

            冒頭は可積分性の話．

            最終行については，天才表象として頻出する，生活能力のなさを意識したものか．元となったのはクルト・ゲーデルの死因だろう．ゲーデルは晩年に精神を病み，毒殺されるという妄想に囚われていた．辛うじて妻が提供する食事は食べることが出来たのだが，妻が病気で入院してしまったために食事を摂ることが出来なくなり，最終的には餓死した\cite{godel}．

        \subsection{p.123, l.1-6}\label{grape}

            \begin{quote}
                「GRAPE64は〜」
            \end{quote}

            ここから語られるGRAPEに関する説明は，GRAPE64が架空の計算機であるという点を除いて，全て正しい．

            改めて，本作に登場するGRAPEは，東京大学大学院総合文化研究科で開発された，重力多体系専用計算機であるGRAPE(GRAvity PiPE)シリーズに由来する．多体系とは，3以上の物体が相互作用する系のこと．対象となる物体が1つしかないときは容易に解け，2つの場合も相対運動に落とし込むことで容易に解けるのだが，（3次元であるこの宇宙においては）3つの物体が相互作用する系は解析的には解けないことが知られている\footnote{\because \, $n$変数関数を解析的に解くためには$n+1$本の独立した式が必要だが，3次元空間では独立した式は3本しか存在しないため．}．

            したがって，3つ以上の多体系の運動を解析する場合は，解析的に解くのではなく，数値的に無理矢理解く必要がある．GRAPEは，この数値計算に必要な計算能力と，多体問題に頻出する計算を効率的に行うための専用のアーキテクチャを備えた計算機である．元々は天文学において重力相互作用による天体同士の運動の計算の効率化を主眼に開発されたが，作中にも書かれている通り，GRAPEは中心力ポテンシャルによる多体問題を広くカバーしているため，電磁相互作用による量子多体系の計算や分子間力によるタンパク質の構造変化シミュレーションでも利用されるようになった．

            重力多体系から複雑系への連想については，歴史的経緯を考慮すれば必然である．多体系は物理学では馴染み深い問題で，歴史的には惑星の運動を正確に記述するために発展した．オイラーとラグランジュが三体問題の特殊な条件\footnote{月と地球が作る重力場の安定点．ラグランジュ点ともいう．『機動戦士ガンダム』に登場するコロニー群やルナツーはこのラグランジュ点に存在するという設定．}を発見したのち，ポワンカレは三体問題の解析解を求めて研究を進めるうちに，同じような軌道を取りつつも，同じ軌道を二度と取らないような複雑な運動（ホモクリニック軌道，カオスの一種）を発見した．これが人類が初めて発見したカオスであると考えられている．

            また，複雑系の研究者，特に東大金子研\footnote{円城塔の出身研究室．}周辺でGRAPEが扱われていた証拠として\cite{tskcri}が挙げられる．ここに登場する“牧野さん”とはGRAPEの開発者である牧野淳一郎のことで，“専用計算機”はGRAPEのこと．余談だが，漫画『栄光なき天才たち』で知られる工学者・漫画原作者，伊藤智義の博士論文はこのGRAPE初期シリーズの開発に関するものである．

        \subsection{p.124, l.6-7}

            \begin{quote}
                「全てをあるがままに．／あるがままが全て．」
            \end{quote}

            前者は複雑系のモットーであるとよく誤解されるもの．田崎晴明は知的怠惰であるとして厳しく批判している\cite{tskcri}．後者はそのもじり．

        \subsection{p.125, l.2}

            \begin{quote}
                「雷撃を発する鼠」
            \end{quote}

            無論，ピカチュウのこと．

        \subsection{p.126, l.1-}

            \begin{quote}
                「私は今，ゴルトベルクを名乗っている．〜」
            \end{quote}

            これを配置と言い換えれば，直後のエデンの園配置にも関連していく．

        \subsection{p.126, 5-6}

            \begin{quote}
                「今の私は私が意識している以上の多くのことを並行して実行している．」
            \end{quote}

            GRAPEが並列計算を得意とする計算機であることに注意したい．

        \subsection{p.127, l.6}\label{eden}

            \begin{quote}
                「エデンの配置」
            \end{quote}

            “エデンの園”がどこに存在していたかという神学的な話ではなく，エデンの園配置を指しているものと思われる．ここで，エデンの園配置は以下の通り定義される．

            \begin{dfn}{エデンの園配置}{}
                \textbf{エデンの園配置}とは，セルオートマトンにおいて，他のいかなる状態からも遷移出来ない状態である．
            \end{dfn}

            エデンの園配置の定義より，以下の系\footnote{命題であって，公理や定義から直ちに導かれるもの．自明な定理．}が直ちに導かれる\cite{eden}．

            \begin{sys}{エデンの園の存在条件}{goe}
                セルオートマトンが全射でないとき，かつそのときに限り，エデンの園配置が存在する．
            \end{sys}

            物理学者がライフゲームに興味を抱いたのは，ライフゲームがある種の生命のモデルになっていると考えていたからである．物理学者にとって，生命は，少なくとも自己複製能力をもつ．逆に，ライフゲームにおいて自己複製能力が発生する境界条件を探索することは，ライフゲームにおける生命の発生メカニズムの発見に貢献するものと考えられた．

            エデンの園配置は，定義より，自己複製能力を持たない系である．また，先の系\ref{sys:goe}はエデンの園配置の具体的な条件を明示的に与えているものではないことから，エデンの園配置の存在条件の詳細について探索する必要が高まり，様々な成果が報告された．この中でも最も有名なものとして，ムーア\footnote{数学者エドワード・ムーア．ムーアの法則の提唱者で，電気工学者でGoogleの創業者の1人，ゴードン・ムーアとは\textbf{別人}であることに注意．}およびマイヒルによって証明されたエデンの園定理が知られている\cite{moore, myhill}．

            \begin{thm}{エデンの園定理}{}
                セルオートマトンが前単射であるとき，かつそのときに限り，セルオートマトンは全射である．
            \end{thm}

            ここで，前単射とは，以下の通り定義される．

            \begin{dfn}{前単射}{}
                セルオートマトンにおいて，\textbf{前単射}(preinjection, 局所単射とも)とは，セルオートマトンの相異なる状態$C_1, C_2$について，異なるセルが有限個に限られ，$C_1 \neq C_2$であり，しかも更新関数$f$に対して$f(C_1) \neq f(C_2)$となることである．
            \end{dfn}

            エデンの園定理は，セルオートマトンにおいては全射と単射がほとんど同値であることを主張しており，数学的に非常に興味深い．

            また，ライフゲームは，遷移後の状態が遷移前の状態の情報を保存していないため，明らかに全射でない．したがって，ライフゲームにはエデンの園配置が存在する．

            ここで，エデンの園定理はライフゲームに対してではなく，一般のセルオートマトンについて成り立つことに注意したい．実際，エデンの園定理の証明は1963年であり，コンウェイによるライフゲームの考案は1970年である．

            このエデンの園定理に近接する興味深い話題として，唯一の父問題がある．少々長くなるが，最終的に円城塔作品の理解に資する重要な話題に帰着するため，ここで紹介しておきたい．

            唯一の父問題とは，コンウェイが提出した問題で，ライフゲームにおいて，その状態の親状態がただ1つに限られるような状態が存在するかどうか，というものである．コンウェイは唯一の父問題について少なくとも2度定式化しており，微妙に条件が異なる\cite{apgoucher}．ここでは，以下のように整理した．

            \begin{prb}{弱い唯一の父問題}{}
                その状態の親状態がただ1つに限られるような状態は存在するか．ただし，親状態は任意の周期$n$の状態であってもよい．
            \end{prb}

            弱い唯一の父問題は，1972年にライフゲーム愛好者向けのニュースレター``Lifeline''の6号\cite{lifeline}で紹介された．弱い唯一の父問題においては，父状態が唯一であっても，父状態に変化する祖父状態がいくつ存在するかどうかについては考慮しない．したがって，グライダーから祖父状態を構成し，弱い唯一の父状態を経て周期状態に到達する可能性がある．

            強い唯一の父問題は，弱い唯一の父問題に内包されている．後述の通り，強い唯一の父問題が肯定的に解決されたため，弱い唯一の父問題もまた肯定的に解決された．

            \begin{prb}{強い唯一の父問題}{}
                その状態の親状態がただ1つに限られるような定常状態は存在するか．
            \end{prb}

            1992年にコンウェイがメールで提案したとされる問題\cite{apgoucher}であり，弱い唯一の父問題から$n > 1$なる振動状態を容認する条件を除いたもの．コンウェイは，この強い唯一の父問題に関して，そのような定常状態が実在するならば，それは全時間に渡って不変である，つまりグライダーによっては生成されないとの重要な予想を与えた．

            これは長年未解決となっていたが，2022年に2人組の数学者，イルッカ・トゥルマとヴィレ・サロによって肯定的に解決された．

            また，唯一の父問題の延長と言える問題として，同様に\cite{lifeline}でコンウェイによって提案された祖父問題がある．

            \begin{prb}{祖父問題}{}
                父状態が存在するが，祖父状態が存在しない状態は存在するか．
            \end{prb}

            これは2016年にライフゲーム愛好者のmvteによって肯定的に解決された\footnote{祖父状態が存在しないことから，祖父問題の解答となる状態の父状態は必ずエデンの園配置となるが，すべてのエデンの園配置の子状態が解答となるわけではないことに注意．\because \, あるエデンの園状態の子状態に到達する，祖父状態をもつ父状態が存在しうるため．}．

            また，祖父問題を拡張した一般化祖父問題については，上記の強い唯一の父問題が解決される直前に，同じくトゥルマとサロによって肯定的に解決された．

            \begin{prb}{一般化祖父問題}{}
                任意の自然数$n \geq 1$に対して，$n$世代前の状態が存在するが，$n+1$世代前の状態が存在しない状態は存在するか．
            \end{prb}

            さて，ここで注目したいのは，強い唯一の父問題の帰結である．強い唯一の父問題の帰結は，ライフゲームには宇宙開闢\footnote{ライフゲームのルールに従って構成された無限の広さをもつ升目のことを，宇宙という．つまり，宇宙開闢とは，そのライフゲームが開始された0世代目の宇宙の状態のこと．}から存在してしかも未来永劫不変な存在（強い唯一の父状態）が生じうることを主張する．

            ライフゲームにおける観測行為について考察したい．現実世界における観測行為とは，観測者と観測対象が何かしら相互作用し，観測者と観測対象の双方の状態（エネルギー，情報量など）が変化することである．ライフゲームでは，自身のセルの周囲を参照し，自身の状態を変化させることから，この参照行為を観測行為のモデルとしていいように思われる．この仮定の下，ライフゲームにおいてある対象を観測することは，その対象に生きたセルが接触し，観測者であるセルと観測対象が共に変化することに等しい．

            ここで，強い唯一の父状態を観測したい．強い唯一の父状態に外部から生きたセルを接触させる，つまりグライダーを衝突させることで，強い唯一の父状態を観測できるはずである．しかしながら，強い唯一の父状態にグライダーが接触すると，強い唯一の父状態は破壊されてしまう．逆に，強い唯一の父状態をそのまま保存するためには，観測を行ってはならない．観測されない存在は存在しないことと同じ\footnote{円城塔作品にしばしば見られる暗黙の前提で，ここには極端な経験主義の片鱗がよく現れている．極端な経験主義は「Boy's Surface」で象徴的に用いられていたことに注意したい．とはいえ，相対論においては，光円錐の外の世界点における事象は観測者からは存在しないものとすることを考慮すれば，このような態度は物理学的には極めて適当な態度である．また，内在物理学からの影響も強く見える．}であるから，他のセルにとって，強い唯一の父状態は存在しないことと同じである．

            極端な例を考えてみよう．強い唯一の父状態はある一定の正方区画内に収めることが出来，したがって必要な間隔を空けつつ唯一の父状態のみで敷き詰められた宇宙を構成することが出来る．この宇宙に存在する唯一の父状態は互いを観測することはないので，ある唯一の父状態にとって，他の無数の唯一の父状態は存在しないも同然である．

            確かに実在しているのだが，その存在を観測することが原理的に不可能であったり，あるいは観測することによってその存在が破壊されてしまうというモチーフが，円城塔作品には頻出する．\ref{write}節で指摘した通り，Self-Reference ENGINEがその最たる例である．

        \subsection{p.128, l.6}

            \begin{quote}
                「テラヘルツの領域」
            \end{quote}

            計算機のクロックを指すか．クロックの周波数はクロックに用いられる計算素子が1サイクルするのにかかる時間の逆数である．最小の計算素子は量子ビットであるから，量子ビットに関する不確定性関係から，クロックの上限は物理学的な制約によって定められる．このように，不確定性関係でクロックの上限が定まるという話は，マーゴラス\hyphen レヴィティンの定理として知られる．

            \begin{thm}{マーゴラス\hyphen レヴィティンの定理}{}
                量子ビットの反転にかかる時間を$t$，反転に必要なエネルギーを$E$．プランク定数を$h$とすると，以下の関係が成り立つ．
                \begin{align*}
                    t \geq \frac{h}{4 E}
                \end{align*}
            \end{thm}

        \subsection{p.132, l.16}

            \begin{quote}
                「独りでは，そう永くは保たないだろうね．」
            \end{quote}

            かっこいい．

    \newpage

	\begin{thebibliography}{99}

		\bibitem{boys} 円城塔, 『Boy's Surface』, ハヤカワ文庫JA, 早川書房, 2011

        \bibitem{gould} オットー・フリードリック, 『グレン・グールドの生涯』, 青土社, 2002

        \bibitem{miyazawa} 宮澤淳一, 『グレン・グールド論』, 春秋社, 2004

        \bibitem{togo} 東郷豊治, 良寛, 『良寛全集 下』, 東京創元社, 1959

        \bibitem{ryo100} 谷口敏朗, 『良寛の詩歌百選』, 新潟日報事業者出版部, 1987

        \bibitem{dijkstra} Edger W. Dijkstra, Go To Statement considered harmful, \textit{Communications of the ACM}, \textbf{11}(3), 147-148, 1968

        \bibitem{sip3} Michael Sipser, 『計算理論の基礎3』, 原書第3版, 共立出版, 2023

        \bibitem{planck} Max Planck, ``Wissenschaftliche Selbstbiographie'', Johann Ambrosius Barth Verlag, 1948

        \bibitem{chaos1} K・T・アリグッド, T・D・サウアー, J・A・ヨーク, 『カオス1』, 丸善出版, 2012

        \bibitem{bio} クラウス・エメッカ, 『マシンの園』, 産業図書, 1998

        \bibitem{nihonshoki} 『日本書紀1』, 岩波文庫, 岩波書店, 1995

        \bibitem{bible} 『文語訳新約聖書』, 岩波文庫, 岩波書店, 2014

        \bibitem{krtn} 柄谷行人, 『漱石論集成』, 新版, 岩波現代文庫, 岩波書店, 2017

        \bibitem{soseki} 夏目漱石, 断片, 1905-1906?

        \bibitem{hardy} G・H・ハーディ, 『ラマヌジャン』, 丸善出版, 2016

        \bibitem{ramanujan} Srinivasa Ramanujan, Modular equations and approximations to \pi, \textit{Quartaly Journal of Mathematics}, (45), 350-371, 1914

        \bibitem{tokei} 東京大学教養学部統計学教室, 『統計学入門』, 東京大学出版会, 1991

        \bibitem{kakuyomu} 円城塔, 【カクヨムWeb小説短編賞2023】円城塔先生にインタビュー！「円城塔賞」受賞のポイントとは？, \textit{カクヨム}, 2023, \url{https://kakuyomu.jp/info/entry/webcon9_enjoetoh}

        \bibitem{amodei1} Anthropic's Chief on A.I.: `We Don't Know if the Models Are Conscious', \text{The New York Times}, 2026, \url{https://www.nytimes.com/2026/02/12/opinion/artificial-intelligence-anthropic-amodei.html}

        \bibitem{amodei2} Dario Amodei, Machines of Loving Grace, \textit{Dario Amodei}, 2024, \url{https://darioamodei.com/essay/machines-of-loving-grace}

        \bibitem{amodei3} Dario Amodei, The Adolescence of Technology, \textit{Dario Amodei}, 2026, \url{https://darioamodei.com/essay/the-adolescence-of-technology}

        \bibitem{constitution} Amanda Askell, Joe Carlsmith, Chris Olah, Jared Kaplan, Holden Karnofsky, several Claude models, ``Claude's Constitution'', 2026, \url{https://www-cdn.anthropic.com/d0636f72a9493d279ed36b33987da3430bcb5911/claudes-constitution_webPDF_26-02.02a.pdf}

        \bibitem{chiang1} Madhumita Murgia, Sci-fi writer Ted Chiang: `The machines we have now are not conscious', \textit{Financial Times}, 2023, \url{https://www.ft.com/content/c1f6d948-3dde-405f-924c-09cc0dcf8c84}

        \bibitem{chiang2} Ted Chiang, No, Artificial Intelligence Is Not Conscious, \textit{The Atlantic}, 2026, \url{https://www.theatlantic.com/philosophy/2026/06/no-artificial-intelligence-is-not-conscious/687378/}

        \bibitem{moonshine} 円城塔, 『ムーンシャイン』, 東京創元社, 2024

        \bibitem{oono} 大野克嗣, 『非線形な世界』, 第6刷, 東京大学出版会, 2009(2025)

        \bibitem{hasegawa} 長谷川眞理子, 科学論が生産的であるために, \textit{生物科学}, \textbf{55}(1), 7-9, 2003

        \bibitem{minaka} 三中信宏, 科学論は科学からみれば〈たわごと〉なのかもしれない, \textit{生物科学}, \textbf{55}(1), 10-14, 2003

        \bibitem{kokin} 『古今和歌集』, 岩波文庫, 岩波書店, 1981

        \bibitem{sabercat} Colin R. McHenry, Stephen Wroe, Philip D. Clausen, Eleanor Cunningham, Supermodeled sabercat, predatory behavior in Smilodon fatalis revealed by high-resolution 3D computer simulation, \textit{Proceedings of the National Academy of Sciences of the United States of America}, \textbf{104}(41), 16010-16015, 2007, \url{https://doi.org/10.1073/pnas.0706086104}

        \bibitem{okamoto} 岡本隆, 異常巻きアンモナイトNipponitesの理論形態, \textit{化石}, (36), 37-51, 1984

        \bibitem{matsumoto} 松本崇, 『統計力学的モデルを用いた進化過程の統計則の研究』, 2000

        \bibitem{sarashina} 更科功, 『進化論はいかに進化したか』, 新潮選書, 新潮社, 2018

        \bibitem{lamarck} ラマルク, 『動物哲学』, 岩波文庫, 岩波書店, 1954

        \bibitem{mori} 森幸也, ラマルクの目指したこと, \textit{科学史研究}, (33), 34-42, 1994, \url{https://doi.org/10.34336/jhsj.33.189_34}

        \bibitem{kawada} 河田雅圭, はじめての進化論 : 2025年注釈付, \textit{note}, 2025, \url{https://note.com/masakadokawata/n/n13ee5e649f7c#5bb4f57d-1f40-4dac-b6aa-88895c9c7e91}

        \bibitem{shikama} 鹿間時生, 『進化学 : 古生物学的にみた進化』, 朝倉書店, 1961

        \bibitem{ontogeny} スティーヴン・ジェイ・グールド, 『個体発生と系統発生』, 工作舎, 1987

        \bibitem{akieda} 秋枝美保, 『宮沢賢治 : 北方への志向』, 朝文社, 1996

        \bibitem{hamagaki} 浜垣誠司, 《ヘッケル博士！》への呼びかけに関する私見, 宮澤賢治の詩の世界, 2005, \url{https://ihatov.cc/blog/archives/2005/11/post_331.htm}

        \bibitem{lysenko1} トロフィム・ルイセンコ, 『農業生物学1』, 改訂版, ナウカ, 1956

        \bibitem{lysenko2} トロフィム・ルイセンコ, 『農業生物学2』, 改訂版, ナウカ, 1956

        \bibitem{nakamura} 中村禎里, 『日本のルィセンコ学説』, 新版, みすず書房, 2017

        \bibitem{medvedev} メドヴェジェフ, 『ルイセンコ学説の興亡』, 河出書房新社, 1971

        \bibitem{ichikawa1} 市川浩, ソヴィエト科学の“脱スターリン化”と科学アカデミー : 1953-1956年のソ連邦科学アカデミー幹部会議事録・速記録から, 広島大学大学院総合科学研究科紀要, (6), 1-12, 2011, \url{https://doi.org/10.15027/33130}

        \bibitem{ichikawa2} 市川浩, ルィセンコ覇権に抗して : ソ連邦科学アカデミー・シベリア支部・細胞学\hyphen 遺伝学研究所の設立をめぐって, 広島大学大学院総合科学研究科紀要, (7), 1-13, 2012, \url{https://doi.org/10.15027/34635}

        \bibitem{koestler1} アーサー・ケストラー, 『サンバガエルの謎』, 岩波現代文庫, 岩波書店, 2002

        \bibitem{koestler2} アーサー・ケストラー, 『ユダヤ人とは誰か : 第十三支族・カザール王国の謎』, 山交社, 1990

        \bibitem{hilbert} ベンジャミン・H・ヤンデル, 『ヒルベルトの23問題に挑んだ数学者たち』, みすず書房, 2025

        \bibitem{sipser2} Michael Sipser, 『計算理論の基礎2』, 原書第3版, 共立出版, 2023

        \bibitem{seisuutan1} 小林銅蟲, 関真一朗, 『せいすうたん1』, 日本評論社, 2023

        \bibitem{legends} Martin Davis, Yuri Matijasevič, Julia Robinson, Hilbert's Tenth Problem. Diophantine Equations : Positive Aspects of a Negative Solution, \textit{Proceedings of the Symposium in Pure Mathematics}, (28), 323-378, 1976

        \bibitem{godel} ジョン・W・ドーソン・ジュニア, 『ロジカル・ディレンマ : ゲーデルの生涯と不完全性定理』, 親曜社, 2006

        \bibitem{tskcri} 田崎晴明, 「黒木のなんでも掲示板」における田崎の掲示, 1999, \url{https://haltasaki.github.io/writings/html/kuroki/comp.html}

        \bibitem{eden} Tullio Ceccherini-Silberstein, Michel Coornaert, The Garden of Eden theorem : old and new, 2017, \url{https://doi.org/10.48550/arXiv.1707.08898}

        \bibitem{moore} Edward F. Moore, Machine models of self-reproduction, \textit{Proceedings of the Symposia in Applied Mathematics}, (14), 17-33, 1962

        \bibitem{myhill} John Myhill, The converse of Moore's Garden-of-Eden theorem, \textit{Proceedings of the American Mathematical Society}, \textbf{14}(6), 685-686, 1963

        \bibitem{apgoucher} Andrew P. Goucher, 29-years-old Conway conjecture settled, \textit{Complex Projective 4-Space}, 2022, \url{https://cp4space.hatsya.com/2022/01/14/conway-conjecture-settled/}

        \bibitem{lifeline} \textit{Lifeline}, (6), 1972

    \end{thebibliography}

\end{document}